% Source: https://artofproblemsolving.com/community/c3h1593863_pls_explain_this_problem_pls
% Stable identity: aops-topic-1593863-q01
In a parallelogram opposite sides are equal, $AB=CD$ and $AD=BC$, so the perimeter is $2(AB+AD)=10$, giving $AB+AD=5$.

\begin{asy}
usepackage("amsmath");
size(4.5cm);
defaultpen(fontsize(10pt));
pair A=(0,0),B=(2.2,0),D=(0.9,1.6),C=B+D;
draw(A--B--C--D--cycle,linewidth(0.9));draw(B--D,blue+linewidth(1));
label("$A$",A,SW);label("$B$",B,SE);label("$C$",C,NE);label("$D$",D,NW);
label("\mbox{}",(min(currentpicture,user=true)+max(currentpicture,user=true))/2,black); // keeps the text after the figure black when LaTeX typesets the labels inline
\end{asy}

The perimeter of triangle $ABD$ is $AB+AD+BD=8$, so
\[
BD=8-(AB+AD)=8-5=\boxed{3\text{ cm}}.
\]
